\chapter{Metodologia di confronto}
\label{chap:metodologia}

Questo capitolo descrive come i tre sensori sono stati confrontati.

La Sezione~\ref{sec:metodo-approccio} distingue le due attività su cui poggia il confronto,
quella sperimentale e quella numerica. La Sezione~\ref{sec:protocollo} fissa il protocollo
con cui sono stati acquisiti i 444 trial del Capitolo~\ref{chap:confronto}, mentre la
Sezione~\ref{sec:metriche} definisce le metriche e le convenzioni con cui sono stati letti.
La Sezione~\ref{sec:protocollo-ld2420} dichiara le deroghe al protocollo adottate per il
secondo radar, e la Sezione~\ref{sec:webui} presenta l'interfaccia web realizzata a supporto
della campagna.

\section{Due strade complementari}
\label{sec:metodo-approccio}

Il confronto poggia su due attività distinte che si alimentano a vicenda.
\begin{itemize}
  \item \textbf{Sperimentale:} consiste nel disporre i sensori e il soggetto in geometrie
    definite e nel cambiare, da una prova all'altra, una sola variabile per volta, che può
    essere la distanza, la quantità o il tipo di movimento, l'angolo oppure il materiale
    interposto. Ogni configurazione viene ripetuta e ogni ripetizione produce un file.
  \item \textbf{Numerica:} i file formano un dataset omogeneo, raccolto nel formato unico
    della Sezione~\ref{sec:c3-4-3} e analizzato a posteriori dalla catena descritta nella
    Sezione~\ref{sec:c3-4-5}. Le metriche, quindi, non vengono lette durante la prova. Sono
    calcolate in seguito, dagli stessi script e con le stesse convenzioni per tutti i
    sensori. In questo modo un risultato si può ricalcolare e una figura si può rifare
    senza ripetere l'acquisizione, e ogni numero del Capitolo~\ref{chap:confronto} risale a
    un file e a una riga di comando precisi.
\end{itemize}

Durante la prova, quindi, l'operatore non ha bisogno di leggere i numeri, e l'interfaccia web serve a controllare in tempo reale che i sensori vedano ciò che devono vedere (Sezione~\ref{sec:webui-ruolo}).

\section{Protocollo sperimentale}
\label{sec:protocollo}
\label{sec:protocollo-configurazione}

Il protocollo è stato definito prima delle acquisizioni. Il PIR lavora nella configurazione a lui più favorevole, con il ponticello di ritrigger in posizione H, in tutte le serie in cui il ponticello conta (Sezione~\ref{sec:interpretazione-pir}), mentre
l'\ldb\ mantiene sempre le soglie di fabbrica (Sezione~\ref{sec:c3-6}). L'\ldc\ è stato
provato in un secondo tempo nelle stesse geometrie, con le deroghe dichiarate nella
Sezione~\ref{sec:protocollo-ld2420}. Ogni scenario è ripetuto più volte, e i valori
riportati nella forma $\mu \pm \sigma$ sono la media e la deviazione standard su quelle
ripetizioni.

\subsection{Setup}
\label{sec:protocollo-setup}

L'\ldb\ e l'\pir\ sono letti dallo stesso firmware, sulla stessa base dei tempi
(Sezione~\ref{sec:c3-4-2}). I due sensori sono fissati uno accanto all'altro su un
supporto di cartone appoggiato sul bordo di un tavolo, a circa 85~cm di altezza
(Figura~\ref{fig:setup-foto}), e puntano verso un'area di prova che ricade nel campo di
entrambi. Nella stanza non ci sono oggetti in movimento (nessun ventilatore, tende ferme,
porta chiusa). Le quattro geometrie di prova sono schematizzate in
Figura~\ref{fig:geometrie}.

\begin{figure}[!htbp]
\centering
\includegraphics[width=\linewidth]{Setup}
\caption{Il banco di prova. (a) Il PIR \pir\ e l'\ldb\ fissati sul supporto di cartone,
con l'ESP32 sulla breadboard accanto. (b) Lo stesso supporto con l'\ldc\ della campagna
ridotta.}
\label{fig:setup-foto}
\end{figure}

\begin{figure}[!htbp]
\centering
\includegraphics[width=\linewidth]{Geometrie}
\caption{Le geometrie della campagna, in scala. (a) La stanza, con i sensori sul tavolo e
le tacche a 1--5~m sull'asse, usata per distanza, dose-risposta, latenze, ostacoli e due
persone. (b) Lo scenario del progetto, con i sensori fissati sotto il piano e la persona
rannicchiata a 50--60~cm. (c) L'arco a 1~m per la copertura angolare. (d) Il corridoio
della prova di portata, con i due vani porta a 4 e 5~m attraverso cui si vedono i punti
più lontani.}
\label{fig:geometrie}
\end{figure}

\subsection{Scenari}
\label{sec:protocollo-scenari}

La Tabella~\ref{tab:scenari} raggruppa gli scenari per domanda. Di norma uno scenario è ripetuto
cinque volte, dieci per la latenza d'ingresso e tre quando serve una seconda persona o le
condizioni da coprire sono molte. Ogni scenario ha un nome, che è anche il nome dei suoi
file. Il respiro a ritmo imposto usa lo stesso setup e lo stesso formato dei file ed è trattato
nel Capitolo~\ref{chap:vitalita}.

\begin{table}[!htbp]
\centering
\caption{Gli scenari della campagna, con i trial per condizione, lo scarto iniziale
applicato da tutti gli script (Sezione~\ref{sec:metriche-transitorio}) e la sezione dei
risultati. La campagna ridotta sull'\ldc\ ripete entro 2~m gli scenari con l'asterisco
(Sezione~\ref{sec:protocollo-ld2420}), con uno scarto di 90~s nei trial da fermo.}
\label{tab:scenari}
\small
\begin{tabularx}{\linewidth}{@{}>{\raggedright\arraybackslash}p{0.27\linewidth}>{\raggedright\arraybackslash}X>{\raggedright\arraybackslash}p{0.13\linewidth}c>{\raggedright\arraybackslash}p{0.13\linewidth}@{}}
\toprule
\textbf{Scenario} & \textbf{Domanda} & \textbf{Trial} & \textbf{Scarto} & \textbf{Risultati} \\
\midrule
Stanza vuota, notte e giorno* & falsi positivi & 7,03~h & 60~s & Sezione~\ref{sec:falsi-positivi} \\
Persona ferma seduta a 2,3~m & falsi negativi & 5 & 30~s & Sezione~\ref{sec:persona-ferma} \\
Dose-risposta a 1~m (immobile, micro-movimenti, cammino sul posto)* & quantità di movimento & 5 per condizione & 20~s & Sezione~\ref{sec:dose-risposta} \\
Cammino sul posto a 1--5~m* & accuratezza della distanza, PIR sul posto & 5 per distanza & 20~s & Sezioni~\ref{sec:distanza} e~\ref{sec:dose-risposta} \\
Attraversamento a 1--5~m & tipo di movimento & 3 per distanza & 20~s & Sezione~\ref{sec:attraversamento} \\
Sotto il banco a 60~cm (immobile, micro-movimenti)* & scenario del progetto & 5 per condizione & 40~s & Sezione~\ref{sec:sotto-banco} \\
Ingresso e uscita dalla stanza* & latenze & 10 e 5 & evento a 30~s & Sezione~\ref{sec:latenza} \\
Arco a 1~m, sei azimut* & copertura angolare & 1--3 per azimut & 20~s & Sezione~\ref{sec:angolare} \\
Selettività con gate massimo 2* & vicino di banco & 3 per scenario & 120~s & Sezione~\ref{sec:angolare} \\
Due persone, tre geometrie* & separazione dei bersagli & 3 per geometria & 20~s & Sezione~\ref{sec:due-persone} \\
Pannelli interposti, a 3~m e a 1~m* & attenuazione & 3 per materiale & 20~s & Sezione~\ref{sec:ostacoli} \\
Corridoio, 5--8~m* & portata massima & 3--5 per distanza & 20~s & Sezione~\ref{sec:portata-massima} \\
\bottomrule
\end{tabularx}
\end{table}

\subsection{Ground truth}
\label{sec:protocollo-ripetizioni}
\label{sec:protocollo-groundtruth}

La presenza e lo stato reali del soggetto sono dichiarati dall'operatore nei metadati di
ogni acquisizione (Sezione~\ref{sec:c3-4-3}) e valgono per l'intero file. Per le latenze
questo non basta. Il soggetto entra o esce quindi a un istante annunciato dallo script di
acquisizione (Sezione~\ref{sec:c3-4-4}), e la latenza si misura a partire da
quell'istante.

\subsection{Registro}
\label{sec:registro}

Ogni sessione è annotata in un registro con data, ora, temperatura della stanza,
soggetto, alimentazione e note. Per ogni scenario, oltre alla distanza, sono annotati
anche il tipo di movimento e la postura, perché il PIR risponde a questi e non alla
distanza (Sezione~\ref{subsec:fresnel}). Le sessioni si sono svolte fra 25 e 29\,\textdegree C. Il registro
completo è nel repository del lavoro~\cite{repoTesi}.

\section{Metriche e convenzioni}
\label{sec:metriche}

\subsection{Metriche}
\label{sec:metriche-elenco}

Le metriche sono calcolate su ogni trial dallo script di analisi della
Sezione~\ref{sec:c3-4-5} e poi aggregate per scenario.
\begin{itemize}
  \item \textbf{Falsi negativi (\%):} campioni con presenza reale in cui il sensore riporta
    assenza. È la metrica dello scenario con la persona ferma, e il suo complemento, cioè
    la frazione di campioni con presenza riportata, è il tasso di rilevamento.
  \item \textbf{Falsi positivi (eventi/h):} transizioni da assenza a presenza per ogni ora
    di stanza vuota. Si contano gli eventi e non i campioni, perché un falso positivo di
    dieci secondi non è dieci volte più grave di uno di un secondo.
  \item \textbf{Latenza di rilevamento e di rilascio:} tempo che passa fra l'evento e la
    prima transizione stabile del sensore, all'ingresso e all'uscita. Per confrontare i
    sensori si usa la differenza appaiata sullo stesso trial, in cui il tragitto e il
    tempo di reazione si cancellano.
  \item \textbf{Accuratezza della distanza:} scostamento fra la distanza riportata dal
    radar e quella reale, tenendo separata la dispersione fra un trial e l'altro da quella
    entro il singolo trial. La linearità della risposta si valuta con una regressione lineare fra distanza riportata e reale. Il coefficiente di determinazione $R^2$ indica quanta parte della variazione della distanza riportata è spiegata dalla retta, e vale 1 quando i punti vi stanno esattamente sopra.
  \item \textbf{Energia:} valore normalizzato 0--100 che l'\ldb\ riporta per il bersaglio e
    per ogni gate (Sezione~\ref{sec:c3-5-1}). Fa da grandezza graduata del movimento e da
    testimone dello stimolo, perché due prove con la stessa energia hanno riprodotto lo
    stesso movimento.
\end{itemize}

\subsection{Zero eventi}
\label{sec:metriche-zero}

Quando in una finestra di osservazione non si verifica alcun evento, non si scrive zero
eventi all'ora ma si riporta il limite superiore dato dalla regola del tre, $3/T$ con $T$
in ore, che approssima il limite al 95\,\% di confidenza per un processo di Poisson senza
eventi~\cite{hanley1983rule}. Questo numero dipende dal tempo di osservazione, non dal sensore.

\subsection{Transitorio iniziale}
\label{sec:metriche-transitorio}

I primi secondi di ogni trial vengono scartati, perché contengono il posizionamento del
soggetto e la coda di presenza del radar (Sezione~\ref{sec:c3-5-3}). La durata dello
scarto dipende dallo scenario, è riportata nella Tabella~\ref{tab:scenari} ed è la stessa in
tutti gli script e in tutte le figure. Negli scenari ordinari vale 20~s. È più lunga sotto
il banco, dove il soggetto deve anche posizionarsi, a stanza vuota, dove l'operatore deve
prima uscire, e negli scenari di selettività, dove la coda del radar dopo il
posizionamento arriva a superare i 100~s (Figura~\ref{fig:transitorio} e
Sezione~\ref{sec:angolare}). Nei trial da fermo dell'\ldc, che rilascia in circa 55~s, lo
scarto sale a 90~s.

\begin{figure}[!htbp]
\centering
\includegraphics[width=\linewidth]{Transitorio}
\caption{Il transitorio iniziale nei tre trial con la persona ferma a 1~m a
90\textdegree\ (gate massimo 2). La presenza riportata è attiva dall'inizio del trial, si
spegne una sola volta e non si riaccende, perché è la coda del posizionamento e non un
rilevamento. Con lo scarto ordinario di 20~s entrerebbe nella statistica, con quello di
120~s resta fuori.}
\label{fig:transitorio}
\end{figure}

\subsection{Validità dei file}
\label{sec:metriche-validita}

Un file entra nelle statistiche solo se è un trial. I controlli del fondo, le prove
tecniche e i file marcati come non validi nel registro restano fuori da ogni tabella
grazie alla lista di esclusione della Sezione~\ref{sec:c3-4-5}. Un file non valido non
viene cancellato, ma resta nel dataset con un nome diverso. Il caso tipico è il radar
rimasto muto per un cavo mosso, che si riconosce da centinaia di frame identici.

\section{Protocollo per l'\ldc{}}
\label{sec:protocollo-ld2420}

L'esemplare di \ldc\ in dotazione ha una portata di circa 2~m (Sezione~\ref{sec:c3-2-6}),
e la campagna sul secondo radar si è svolta per intero entro quel limite. Sono stati
ripetuti tutti i test dell'\ldb\ che non richiedono distanze maggiori, per un totale di 161
trial validi, e sono rimaste fuori solo le prove che le richiedono. Entrano invariate le due
più importanti per il progetto, la dose-risposta a 1~m e la persona sotto il banco.

Vanno dichiarate quattro deroghe alla parità di configurazione. I dati sono letti in
energy mode (Sezione~\ref{sec:c3-2-5}), l'unica modalità che esponga le energie per gate.
Le soglie sono tarate sul fondo misurato e il gate massimo è fissato a 6, per le ragioni
della Sezione~\ref{sec:c3-6}. Inoltre, il gate~0 è portato sopra la coda del rumore, nella
configurazione v2 (Sezione~\ref{sec:falsi-positivi}). Infine, il PIR non è cablato nei
test del solo \ldc, e i suoi valori sono marcati come assenti (Sezione~\ref{sec:c3-4-3}).
Poiché l'esemplare in dotazione è con ogni probabilità difettoso, i risultati ottenuti descrivono quel singolo modulo e non vanno
generalizzati al modello \ldc.

Due accorgimenti operativi completano il protocollo. Dopo ogni accensione nessuno resta
entro 2~m dal modulo per 90~s, perché il fondo viene stimato all'avvio e un corpo vicino lo
falserebbe. Prima di ogni sessione, poi, un minuto di stanza vuota serve a verificare che
il fondo per gate sia quello atteso.

\section{L'interfaccia web di supporto}
\label{sec:webui}
\label{chap:webui}

L'ESP32 pubblica i dati e una pagina web li mostra in tempo reale, li salva insieme alle
statistiche di sessione e li esporta in CSV. Questa sezione ne riassume requisiti, progetto, realizzazione e verifica,
e chiarisce il ruolo che ha avuto nella campagna.

\subsection{Requisiti e architettura}
\label{sec:webui-requisiti}

Dall'obbiettivo discendono cinque requisiti funzionali.
\begin{itemize}
  \item \textbf{R1:} l'ESP32 deve rendere disponibili i dati in rete.
  \item \textbf{R2:} la visualizzazione deve seguire il sensore alla cadenza
    dell'acquisizione, cioè 5~Hz.
  \item \textbf{R3:} l'applicazione deve calcolare le statistiche di sessione.
  \item \textbf{R4:} l'applicazione deve esportare i dati in CSV.
  \item \textbf{R5:} l'applicazione deve mostrare l'indice di vitalità calcolato a bordo
    (Capitolo~\ref{chap:vitalita}).
\end{itemize}
A questi si aggiungono due vincoli, che nascono dal contesto.
\begin{itemize}
  \item \textbf{N1:} tutto deve funzionare senza appoggiarsi ad alcuna rete esistente,
    perché in emergenza sismica non c'è connettività, né verso Internet né verso una rete
    locale (Sezione~\ref{sec:dipme}).
  \item \textbf{N2:} il CSV esportato dal browser deve avere lo stesso formato di quello
    prodotto dalla catena seriale, così che in tutta la tesi esista un solo formato dati.
\end{itemize}
La Tabella~\ref{tab:webui-opzioni} riassume le tre architetture considerate.

\begin{table}[!htbp]
\centering
\caption{Opzioni architetturali valutate per la web UI. È stata scelta la A.}
\label{tab:webui-opzioni}
\small
\begin{tabularx}{\linewidth}{@{}p{0.34\linewidth}XX@{}}
\toprule
\textbf{Architettura} & \textbf{Vantaggi} & \textbf{Svantaggi} \\
\midrule
\textbf{A. Sito servito dall'ESP32} (web server e WebSocket a bordo) &
  nessuna infrastruttura, funziona offline, coerente con lo scenario di emergenza &
  RAM e flash limitate, pochi client simultanei \\
B. ESP32, broker MQTT, server con database &
  scalabile, storico illimitato, più vicino a una piattaforma reale &
  richiede broker e server sempre accesi, infrastruttura eccessiva \\
C. Piattaforma cloud (ThingSpeak, Blynk) &
  pronta all'uso &
  non funziona offline, cadenza limitata a circa 15~s, dati fuori controllo \\
\bottomrule
\end{tabularx}
\end{table}

È stata scelta l'opzione~A, perché un nodo che serve la propria dashboard senza alcuna
rete esterna è un caso concreto di scenario senza connessione, cioè proprio lo scenario
del progetto. L'opzione~C viola il vincolo N1, mentre la~B lo rispetta solo con un server
locale che in emergenza non ci sarebbe. Il volume dei dati è piccolo, una trentina di
valori a 5~Hz per circa 2~kB/s, e a bordo non serve persistenza, perché lo storico di
sessione si accumula nel browser ed è il browser a produrre il CSV. L'ESP32 crea quindi una
propria rete WiFi in modalità Access Point e non tenta mai di collegarsi ad altre.
L'opzione~B resta documentata come alternativa valutata e scartata.

\subsection{Progetto e realizzazione}
\label{sec:webui-progetto}
\label{sec:webui-piano}
\label{sec:webui-lezioni}
\label{sec:webui-riferimenti}

\paragraph{Firmware} Il firmware del nodo è scritto in C++ con il framework Arduino, come
il logger della Sezione~\ref{sec:c3-4-2}, di cui riprende per intero la lettura dei sensori
aggiungendo la parte di rete. Il server web e il canale WebSocket sono forniti dalla
libreria ESP Async WebServer, nella versione mantenuta dal progetto
ESP32Async~\cite{esp32async}, perché quella originale non compila più con le versioni
attuali del supporto ESP32 per Arduino. La libreria è asincrona, cioè gestisce le
richieste dei browser in un task separato, eseguito in parallelo al ciclo principale del
programma. In questo modo la lettura del radar prosegue regolarmente a 5~Hz anche mentre
un browser si collega o scarica la pagina. La pagina è incorporata nel firmware come dati
compressi salvati in flash, per circa 80~kB. Per farli stare serve lo schema di partizione
che riserva 3~MB all'applicazione, ed è l'unica impostazione da cambiare rispetto al
logger. Un piccolo server DNS a bordo, infine, risponde con l'indirizzo del nodo a
qualunque nome venga richiesto, così la dashboard si apre da sola appena ci si collega alla
rete del nodo, senza bisogno di conoscerne l'indirizzo IP.

\paragraph{Pagina} La pagina, cioè la parte che gira nel browser, è composta da tre file
(HTML, CSS e JavaScript) per circa 800 righe, con il codice scritto in JavaScript puro,
senza librerie di sviluppo come Angular o React e senza strumenti di compilazione, quindi
i file scritti sono esattamente quelli eseguiti dal browser. Il tempo reale è affidato al
canale WebSocket, disponibile in ogni browser senza librerie, e l'unica dipendenza è la
libreria di grafici Chart.js~\cite{chartjs}, servita dall'ESP32 insieme alla pagina. Come
riferimento di interfaccia è stato esaminato un configuratore web open source per
l'\ldb~\cite{ld2410configurator}, da cui è ripresa la rappresentazione a barre delle
energie per gate con la soglia sovrapposta. La pagina è divisa in quattro aree
(Figure~\ref{fig:dashboard} e~\ref{fig:dashboard-sessione}). La prima
mostra lo stato istantaneo dei due sensori. La seconda contiene la gauge dell'indice di
vitalità, con i colori del semaforo (Sezione~\ref{sec:vitalita-classi}). La terza
raccoglie i tre grafici degli ultimi 60~s, con l'energia, le barre per gate con le soglie
lette dal modulo all'avvio e la timeline della presenza di radar e PIR sulla stessa base
dei tempi. La quarta gestisce la sessione di registrazione, con gli stessi metadati
sperimentali della catena seriale. Le soglie, deliberatamente, non si possono modificare
dalla pagina, perché devono restare fisse per tutta la campagna
(Sezione~\ref{sec:protocollo-configurazione}).

\paragraph{Protocollo dati} Il nodo trasmette i dati ai browser collegati attraverso il
canale WebSocket, con un messaggio per ogni campione, cioè cinque volte al secondo. Ogni
messaggio è un oggetto JSON che contiene tutte le grandezze del CSV del logger più le due
calcolate a bordo (Tabella~\ref{tab:webui-messaggio}). Il nodo conserva inoltre uno storico
di 60~s e lo invia a ogni client che si collega, così il grafico è già pieno anche dopo un
aggiornamento della pagina. Le statistiche sono calcolate nel browser una volta al secondo.
Il CSV esportato riporta, nello stesso ordine, le colonne del logger e i metadati dello
script di acquisizione, più due colonne in coda con l'indice e la classe calcolati dal
firmware (Sezione~\ref{sec:c3-4-3}). Sono proprio questi i dati usati per verificare il
porting (Sezione~\ref{sec:vitalita-bordo}).

\begin{figure}[!htbp]
\centering
\includegraphics[width=\linewidth]{Dashboard}
\caption{La dashboard servita dal nodo, con la testata con lo stato di presenza, lo stato
live dei due sensori (A), l'indice di vitalità calcolato a bordo (B) e i tre grafici degli
ultimi 60~s (C). I dati sono quelli di una sessione registrata dall'ESP32 con una persona
immobile a 1~m, riprodotti a 5~Hz nella stessa pagina. Il radar resta attivo per tutto il
tempo, mentre il PIR non rileva nulla.}
\label{fig:dashboard}
\end{figure}

\begin{figure}[!htbp]
\centering
\includegraphics[width=\linewidth]{Dashboard_sessione}
\caption{La parte inferiore della dashboard in una sessione dal vivo di 60~s, con una
persona davanti al sensore. A sinistra il riquadro D con i metadati della sessione e le
statistiche calcolate nel browser, a destra la diagnostica dell'ESP32 con la build,
l'indirizzo, la memoria libera, i parametri letti dal radar e l'impronta della pagina.}
\label{fig:dashboard-sessione}
\end{figure}

\begin{table}[!htbp]
\centering
\caption{Campi del messaggio WebSocket inviato a ogni client collegato, uno per campione
(5~Hz), con un valore di esempio.}
\label{tab:webui-messaggio}
\small
\begin{tabularx}{\linewidth}{@{}>{\raggedright\arraybackslash}p{0.24\linewidth}>{\raggedright\arraybackslash}p{0.33\linewidth}X@{}}
\toprule
\textbf{Campo} & \textbf{Esempio} & \textbf{Significato} \\
\midrule
\texttt{t} & \texttt{123456} & istante del campione, in millisecondi dall'accensione del nodo \\
\texttt{radar\_ok} & \texttt{1} & 1 se nel periodo è arrivato un frame valido dal radar \\
\texttt{presence}, \texttt{moving}, \texttt{still} & \texttt{1, 1, 0} & presenza complessiva, bersaglio in movimento, bersaglio fermo \\
\texttt{mdist}, \texttt{sdist} & \texttt{145, 0} & distanza del bersaglio in movimento e di quello fermo, in cm \\
\texttt{menergy}, \texttt{senergy} & \texttt{72, 0} & energia del bersaglio in movimento e di quello fermo, 0--100 \\
\texttt{pir} & \texttt{1} & uscita del PIR, letta nello stesso campione \\
\texttt{gates\_m} & \texttt{[0,5,60,10,0,0,0,0,0]} & energia per gate 0--8, canale di movimento \\
\texttt{gates\_s} & \texttt{[0,0,45,8,0,0,0,0,0]} & energia per gate 0--8, canale stazionario \\
\texttt{light}, \texttt{out} & \texttt{24, 1} & livello del fotosensore e stato del pin OUT \\
\texttt{vitality} & \texttt{54} & indice di vitalità calcolato a bordo, 0--100 \\
\texttt{vitality\_class} & \texttt{vitalita\_moderata} & classe dell'indice di vitalità \\
\bottomrule
\end{tabularx}
\end{table}

\paragraph{Realizzazione} Lo sviluppo è stato organizzato in cinque passi, ciascuno con un
criterio di accettazione da verificare prima di passare al successivo.
\begin{enumerate}
  \item \textbf{Rete e pagina statica:} il passo era superato quando la pagina si apriva da
    telefono e da portatile.
  \item \textbf{Lettura dei sensori e WebSocket:} i valori in pagina dovevano seguire il
    sensore senza accumulare latenza.
  \item \textbf{Grafici e storico:} spostando il soggetto di 0,75~m, le barre per gate
    dovevano indicare il gate corretto.
  \item \textbf{Sessione ed esportazione:} il CSV del browser è stato verificato come
    descritto nella Sezione~\ref{sec:webui-verifica}.
  \item \textbf{Indice di vitalità a bordo:} a regime l'indice doveva coincidere entro
    $\pm 1$ con quello calcolato in Python (Sezione~\ref{sec:vitalita-bordo}).
\end{enumerate}
Un problema incontrato lungo la strada merita di essere segnalato a chi ripete il lavoro.
I callback del server asincrono girano su un task diverso dal ciclo principale, e una
stampa di log alla connessione di un client ha corrotto una riga CSV mentre veniva
scritta. Per questo gli eventi di rete passano ora da una coda e vengono stampati dal
ciclo principale fra una riga e l'altra.

\subsection{Verifica di equivalenza}
\label{sec:webui-verifica}

Il vincolo N2 lega l'interfaccia al resto della tesi, e la sua verifica è stata impostata
come un test di equivalenza. La stessa sessione è stata registrata contemporaneamente in
due modi (Figura~\ref{fig:percorsi-dati}), dal browser con la funzione di esportazione e
dalla seriale con lo script di acquisizione lanciato in parallelo sulla porta del nodo. I
due file sono stati poi confrontati campione per campione, allineandoli sul tempo
dell'ESP32. Sui 1346 campioni comuni, cioè sull'intera sovrapposizione temporale, tutte le
colonne coincidono in ogni valore, comprese le energie per gate e le distanze, e lo script
di analisi legge entrambi i file producendo le stesse metriche. Non si tratta quindi di due
formati compatibili, ma dello stesso dato raggiunto per due vie.

\begin{figure}[!htbp]
\centering
\includegraphics[width=\linewidth]{Percorsi_dati}
\caption{I due percorsi dei dati del firmware della web UI. Sopra il percorso seriale, in
cui lo script di acquisizione aggiunge i metadati sperimentali. Sotto la web UI, servita dal nodo in modalità Access Point, che esporta dal browser lo stesso CSV. I due file coincidono
campione per campione e sono letti dallo stesso script di analisi.}
\label{fig:percorsi-dati}
\end{figure}

\subsection{Ruolo nella campagna}
\label{sec:webui-ruolo}

Nella campagna l'interfaccia ha svolto due funzioni.
\begin{itemize}
  \item \textbf{Controllo in tempo reale:} durante il posizionamento dei sensori e del
    soggetto permette di verificare, prima di registrare, che ogni sensore veda ciò che
    deve vedere. Il caso più evidente è il corridoio, dove la timeline della presenza di
    radar e PIR mostra dal vivo il risultato del confronto.
  \item \textbf{Acquisizione della prova di portata:} la prova in corridoio
    (Sezione~\ref{sec:portata-massima}) è stata acquisita con il firmware della web UI, e
    grazie alla verifica di equivalenza quei dati sono entrati nella catena di analisi
    senza alcuna modifica.
\end{itemize}
Il calcolo delle metriche di confronto resta invece affidato alla catena di analisi sui file (Sezione~\ref{sec:metodo-approccio}).
